\documentclass[11pt]{article}

\usepackage{amsmath,amssymb}
\usepackage{enumitem}
\usepackage{geometry}
\geometry{margin=1in}

\title{Sample Oral Exam Questions and Answers\\
Immersed Boundary Method: Discretization}
\author{}
\date{}

\begin{document}
\maketitle

\section*{Overview of IB Discretization}

\begin{enumerate}[label=\arabic*.]

\item \textbf{What are the three main discretization components of the immersed boundary method?}

\textit{Sample answer:}  
The IB method discretizes:
\begin{itemize}
\item the fluid variables on a fixed Eulerian Cartesian grid,
\item the structure on a Lagrangian mesh of material points, and
\item the interaction between them using regularized delta functions.
\end{itemize}

\item \textbf{Describe the basic sequence of operations performed in one IB time step.}

\textit{Sample answer:}  
At each time step, elastic forces are computed on the structure, these forces are spread to the fluid grid, the fluid equations are solved, fluid velocities are interpolated back to the structure, and the structure is advanced in time.

\end{enumerate}

\section*{Eulerian Discretization of the Fluid}

\begin{enumerate}[label=\arabic*., resume]

\item \textbf{How are the fluid variables discretized in space and time?}

\textit{Sample answer:}  
In the classical prototype, velocity and pressure are discretized on a uniform Cartesian grid with spacing $h = \Delta x$, and time is discretized with a step $\Delta t$. Spatial derivatives are commonly approximated with second-order finite differences on a staggered MAC grid. Modern IB solvers may instead use locally refined Cartesian grids and other compatible fluid discretizations. (Historical note: Peskin's original and \emph{Acta Numerica} (2002) schemes use a collocated grid with the centered operator $D^0$, which is why his four-point kernel imposes the even--odd condition; the staggered MAC grid is standard in Roma, Peskin, and Berger (1999) and in IBAMR.)

\item \textbf{Why are staggered grids commonly used in IB implementations?}

\textit{Sample answer:}  
Staggered grids help avoid pressure–velocity decoupling and improve discrete enforcement of incompressibility by placing velocity and pressure at different locations on the grid.

\item \textbf{Write the semi-implicit discrete Navier--Stokes update used in Peskin-style IB methods.}

\textit{Sample answer:}  
A common IMEX prototype is
\[
\rho\left(\frac{\mathbf u^{n+1}-\mathbf u^n}{\Delta t}
+ \mathcal S_h(\mathbf u^n)\mathbf u^n\right)
= -G_h p^{n+1}
+ \mu L_h\mathbf u^{n+1} + \mathbf f^n,
\]
with the incompressibility constraint
\[
D_h\mathbf u^{n+1} = 0.
\]
Here $G_h$, $D_h$, and $L_h$ are compatible discrete gradient, divergence, and vector-Laplacian operators. Equivalently, the pressure-gradient term may be moved to the left with a plus sign.

\item \textbf{Why is this system linear in the unknowns $u^{n+1}$ and $p^{n+1}$?}

\textit{Sample answer:}  
Because the nonlinear convection term is evaluated using $\mathbf u^n$ and, in this explicit-coupling prototype, $\mathbf f^n$ is already known. The unknown velocity and pressure therefore occur only in the linear viscous, pressure, and incompressibility terms. A fully implicit force law would instead lead to a coupled nonlinear problem.

\end{enumerate}

\section*{Lagrangian Discretization of the Structure}

\begin{enumerate}[label=\arabic*., resume]

\item \textbf{How is an immersed structure discretized in the IB method?}

\textit{Sample answer:}  
The configuration $\boldsymbol\chi(q,t)$ is approximated on a Lagrangian mesh or at material sample points. Each interaction point carries the appropriate quadrature weight: approximately $\Delta q$ for a uniformly sampled fiber, a surface-area weight for a membrane, or a reference-volume weight for a solid. In IBFE, structural nodes and interaction quadrature points need not coincide.

\item \textbf{Where are elastic forces computed, and why?}

\textit{Sample answer:}  
Elastic forces are computed entirely in Lagrangian coordinates, where material connectivity and constitutive laws are most naturally defined.

\end{enumerate}

\section*{Force Spreading and Velocity Interpolation}

\begin{enumerate}[label=\arabic*., resume]

\item \textbf{Write the discrete force-spreading formula.}

\textit{Sample answer:}  
\[
f^n(x_i) = \sum_m F^n_m \, \delta_{h}(x_i - X^n_m)\, w_m ,
\]
where $w_m$ is the Lagrangian quadrature weight of interaction point $m$; for a
uniformly parameterized curve, $w_m = \Delta q$.

\item \textbf{Write the discrete velocity interpolation formula.}

\textit{Sample answer:}
The structure velocity is interpolated from the \emph{new} fluid velocity, at
time level $n+1$:
\[
U^{n+1}_m = \sum_i u^{n+1}(x_i)\, \delta_{h}(x_i - X^n_m)\, (\Delta x)^d .
\]

\item \textbf{Why is the same regularized delta function used in both formulas?}

\textit{Sample answer:}  
Using the same kernel is part of the construction, but it is not sufficient by itself. With compatible Eulerian and Lagrangian weights and inner products, interpolation is chosen as the weighted adjoint of spreading. Then
\[
\langle \mathbf u,S_h\mathbf F\rangle_E
=\langle J_h\mathbf u,\mathbf F\rangle_L,
\]
so the two grids exchange equal discrete power.

\end{enumerate}

\section*{Updating the Structure}

\begin{enumerate}[label=\arabic*., resume]

\item \textbf{How is the structure advanced in time in the simplest IB scheme?}

\textit{Sample answer:}  
Using a forward Euler update in which the structure is advected by the velocity
interpolated from the new fluid field:
\[
X^{n+1}_m = X^n_m + \Delta t \, U^{n+1}_m ,
\]
matching the time levels used in decks 1-2 and 1-4.

\item \textbf{What limits the allowable time step in explicit IB schemes?}

\textit{Sample answer:}  
The time step can be limited by explicit advection and by eigenvalues of the discretized elastic coupling. For example, the largest eigenvalue of a fourth-difference bending operator grows like $\mathcal O((\Delta q)^{-4})$ as the material mesh is refined. This is a statement about mesh-scale stiffness, not a claim that a resolved physical bending force must grow like $(\Delta q)^{-4}$.

\end{enumerate}

\section*{Codimension and Structural Forces}

\begin{enumerate}[label=\arabic*., resume]

\item \textbf{What is meant by the codimension perspective in the IB method?}

\textit{Sample answer:}  
Codimension is the difference between the dimensions of the fluid domain and the immersed material domain: for example, a curve in a two-dimensional fluid has codimension one, while a curve in a three-dimensional fluid has codimension two. The same distributional coupling pattern applies, but the material measure, quadrature weights, units of force density, and numerical resolution requirements change with dimension.

\item \textbf{List and briefly describe common force components used for fibers.}

\textit{Sample answer:}  
Typical forces include target-point (penalty) forces that tether points to prescribed locations, stretching forces that resist changes in arc length, and bending forces that penalize curvature.

\item \textbf{Why do bending forces impose especially severe time-step restrictions?}

\textit{Sample answer:}  
Because discrete bending forces involve fourth-order derivatives, the largest eigenvalue of the discrete bending operator scales as $(\Delta q)^{-4}$, leading to very small stable time steps when treated explicitly. As above, this is a statement about mesh-scale stiffness of the discrete operator, not a claim that a resolved physical bending force must grow like $(\Delta q)^{-4}$.

\end{enumerate}

\section*{Weak Enforcement of No-Slip}

\begin{enumerate}[label=\arabic*., resume]

\item \textbf{How is the no-slip condition enforced in the immersed boundary method?}

\textit{Sample answer:}  
No-slip is enforced weakly by interpolating fluid velocity to the structure using the regularized delta function and advancing the structure with that velocity, rather than enforcing velocity equality pointwise.

\item \textbf{What is meant by grid-scale slip, and what happens as $h \to 0$?}

\textit{Sample answer:}  
At finite $h$, the interpolated velocity is an average over the kernel support rather than an exact pointwise trace, so a regularized or grid-scale discrepancy may remain. Under a stable and consistent refinement of the fluid grid, material mesh or quadrature, kernel, and time step, the approximation is expected to converge to the continuum kinematic relation. The limit $h\to0$ alone is not a convergence proof.

\end{enumerate}

\section*{Discrete Delta Functions}

\begin{enumerate}[label=\arabic*., resume]

\item \textbf{Describe the tensor-product construction of a discrete delta function.}

\textit{Sample answer:}  
The discrete delta is built as a tensor product of 1D kernels:
\[
\delta_{\Delta x}(x) = \frac{1}{(\Delta x)^d}\prod_{\alpha=1}^d \phi\left(\frac{x_\alpha}{\Delta x}\right),
\]
where $\phi$ is compactly supported and symmetric.

\item \textbf{What properties of the cosine kernel make it suitable for IB coupling?}

\textit{Sample answer:}  
The four-point-support cosine kernel is symmetric, compactly supported, satisfies the discrete partition-of-unity condition, and has a constant discrete sum of squares. It is simple to implement. Unlike Peskin's classical four-point kernel, however, it does not satisfy the exact first-moment condition for every subgrid offset, so the appropriate kernel depends on the required invariance and accuracy properties. Quantitatively, $M_1(r)=\sum_j(r-j)\phi_{\cos}(r-j)=r-\tfrac12+\tfrac12\bigl(\cos\tfrac{\pi r}{2}-\sin\tfrac{\pi r}{2}\bigr)$ for $0\le r<1$, with maximum $|M_1|\approx0.021$ (in units of $h$) near $r\approx0.21$ and $0.79$. It does satisfy the even--odd condition, and it is the simple approximation to the four-point kernel given by Peskin (2002, eq.~6.28).

\end{enumerate}

\section*{From Classical to Modern IB Methods}

\begin{enumerate}[label=\arabic*., resume]

\item \textbf{What are the main limitations of the classical Peskin (1977), \emph{J. Comput. Phys.} 25:220--252, discretization?}

\textit{Sample answer:}  
In the classical explicit, regularized-delta implementation, important limitations include severe time-step restrictions for stiff structures and reduced accuracy for pressure, traction, or other nonsmooth interfacial quantities. Periodic domains, uniform grids, FFT solvers, and explicit coupling are influential implementation choices, not defining assumptions of the continuum IBM.

\item \textbf{How do modern IB methods address these limitations?}

\textit{Sample answer:}  
Modern methods introduce implicit or semi-implicit coupling, finite-element structural models (IBFE), adaptive mesh refinement, and multilevel solvers to improve stability, accuracy, and flexibility.

\end{enumerate}

\section*{References}
\sloppy

The model answers are adapted from deck 1-4, which follows these sources.
\begin{itemize}
\item C.~S. Peskin, ``Numerical analysis of blood flow in the heart,'' \emph{J. Comput. Phys.} 25 (1977), 220--252, doi:10.1016/0021-9991(77)90100-0.
\item C.~S. Peskin and D.~M. McQueen, ``A three-dimensional computational method for blood flow in the heart. I. Immersed elastic fibers in a viscous incompressible fluid,'' \emph{J. Comput. Phys.} 81 (1989), 372--405.
\item C.~S. Peskin, ``The immersed boundary method,'' \emph{Acta Numerica} 11 (2002), 479--517, doi:10.1017/S0962492902000077. Fiber forces, the IMEX fluid update with $D^0$, kernels (cosine kernel: eq.~6.28).
\item A.~M. Roma, C.~S. Peskin, and M.~J. Berger, ``An adaptive version of the immersed boundary method,'' \emph{J. Comput. Phys.} 153 (1999), 509--534, doi:10.1006/jcph.1999.6293.
\item F.~H. Harlow and J.~E. Welch, ``Numerical calculation of time-dependent viscous incompressible flow of fluid with free surface,'' \emph{Phys. Fluids} 8 (1965), 2182--2189, doi:10.1063/1.1761178 (MAC grid).
\item B.~E. Griffith and N.~A. Patankar, ``Immersed methods for fluid--structure interaction,'' \emph{Annu. Rev. Fluid Mech.} 52 (2020), 421--448, doi:10.1146/annurev-fluid-010719-060228.
\end{itemize}

\end{document}
